% ---------------------------------------------------------------------------
% Improvement: robustness to lighting and rotation
% ELEC 872 - Lab (workshop, Oct 6)
%
% Overleaf: create a new project and paste this file (or upload it).
% Compiles as-is (only needs booktabs).
% ---------------------------------------------------------------------------
\documentclass[11pt]{article}

\usepackage[margin=1in]{geometry}
\usepackage{booktabs}
\usepackage{caption}
\captionsetup{font=small, labelfont=bf}

\title{Improving LBP Face Recognition under\\ Lighting Changes and Rotation}
\author{ELEC 872 --- Lab}
\date{October 6, 2026}

\begin{document}
\maketitle

% ---------------------------------------------------------------------------
\section{What we changed}

Two changes were made, one for each failure mode. Everything else (descriptor,
distance, split) is unchanged, so the comparison against the baseline is fair.

\paragraph{1. Illumination normalisation (lighting).}
A one-side light is a smooth multiplicative field $L(x)$, so the observed image is
$L\cdot(\text{face})$. Blurring the image gives approximately
$L\cdot\text{blur}(\text{face})$; dividing the image by its own Gaussian-blurred copy
therefore cancels $L$ while keeping the local face texture that LBP encodes. This is
implemented in \texttt{lbpface/preprocess.py} as \texttt{normalize\_illumination} and
applied to every image before the LBP step.

\paragraph{2. Rotation-augmented gallery (rotation).}
At enrolment, rotated copies of every gallery image are stored ($-30^\circ$ to
$+30^\circ$ in $5^\circ$ steps). A rotated probe then matches the correctly oriented
template of its own subject instead of failing against upright templates only. This is
implemented in \texttt{experiments/improved.py} as \texttt{fix\_rotation}.

The evaluation lives in \texttt{experiments/improved.py}; the helpers
\texttt{light(img, s)} and \texttt{spin(img, deg)} build the hard probe images.

% ---------------------------------------------------------------------------
\section{Results}

\begin{table}[htbp]
  \centering
  \caption{Rank-1 recognition rate: baseline vs.\ improved.}
  \label{tab:improved}
  \begin{tabular}{lrr}
    \toprule
    Test condition & Baseline & Improved \\
    \midrule
    Normal        & 0.975 & 0.975 \\
    \midrule
    Uniform       & 0.975 & 0.975 \\
    \midrule
    Lighting $s = 0.2$ & 0.960 & 0.975 \\
    Lighting $s = 0.4$ & 0.945 & 0.970 \\
    Lighting $s = 0.6$ & 0.895 & 0.985 \\
    Lighting $s = 0.8$ & 0.705 & 0.975 \\
    Lighting $s = 1.0$ & 0.335 & 0.905 \\
    \midrule
    Rotation $5^\circ$  & 0.970 & 0.985 \\
    Rotation $10^\circ$ & 0.925 & 0.990 \\
    Rotation $15^\circ$ & 0.800 & 0.980 \\
    Rotation $20^\circ$ & 0.630 & 0.980 \\
    Rotation $25^\circ$ & 0.510 & 0.960 \\
    Rotation $30^\circ$ & 0.340 & 0.965 \\
    \bottomrule
  \end{tabular}
\end{table}

The accuracy under the change conditions rises clearly, while the normal-photo
accuracy stays at its baseline value (0.975).

\end{document}
